\section{Discussion}

This study formulates tissue-specific preference in MHC-I presentation as a shared-learning problem over peptide pairs matched by MHC restriction, parent protein, and total positive tissue-occurrence count. Rather than predicting presentation in general, the benchmark asks which member of a controlled pair has stronger evidence in a specified tissue--MHC context. The results show that peptide sequence retains tissue-associated predictive information after these controls are applied. The same task definition is learnable in independently trained human and mouse models, although the species differ in scale and task inventory. Component ablations further support the multi-kernel encoder in both species and the MHC-specific pathway in human, indicating that local sequence patterns and allele-restricted sharing contribute to prediction.

The interpretation analyses connect these predictive effects to specific sequence positions. Pair-matched attribution and exhaustive mutagenesis independently emphasize P2, P3, and P9, while HLA-stratified maps show distinct residue--position structures rather than one universal pooled motif. Changing the tissue task while holding peptide, HLA, and substitution fixed produces greater perturbation dispersion than changing the training seed for every multi-tissue HLA examined. These results demonstrate fitted-model tissue conditioning, but they do not identify which cleavage, transport, trimming, loading, binding, or other tissue-dependent process generated the observed associations.

TissuePMHC addresses a different question from general peptide--MHC predictors. The near-random performance of MHCflurry and NetMHCpan on the matched benchmark indicates that overall binding or presentation scores do not explain the tissue-specific rankings; it does not imply failure at their intended tasks. Similarly, the weak MHC-only control shows that peptide--MHC information without tissue identity cannot satisfy labels that vary across tissues. TissuePMHC may therefore support antigen prioritization and experimental follow-up, but it is not yet suitable for clinical or safety decisions.

Several limitations define the scope of these findings. The comparison peptides are pseudo-negatives based on missing database evidence rather than confirmed non-presented peptides. The source records are observational, and matching occurrence count does not control study, laboratory, donor, platform, detection depth, or uneven tissue and allele coverage. Peptides and parent proteins may also recur across tasks, and the fixed test evaluates represented tissue--MHC combinations rather than unseen tissues, alleles, studies, or cohorts. Finally, attribution, mutagenesis, and same-HLA counterfactuals describe fitted-model dependence rather than biological causality.

Future validation should prioritize protein- or study-disjoint splits, temporal evaluation, independent cohorts, and tissue-label permutation controls. Prospective immunopeptidomics and synthesized substitutions are needed to determine whether the P2/P3/P9 patterns persist outside the present benchmark. Overall, TissuePMHC establishes a sequence-dependent, tissue-conditioned ranking signal in represented MHC-I contexts, while independent and prospective validation remains necessary before biological or clinical application.
